% Instructor copy. Keep separate from the student exam.
% Compile from the course root with pdflatex.
\documentclass[a4paper,11pt,twoside]{article}
\usepackage[T1]{fontenc}
\usepackage{ebgaramond}
\usepackage{amsmath,amssymb}
\usepackage[ebgaramond]{newtxmath}
\usepackage[margin=2cm,includefoot,footskip=8mm]{geometry}
\usepackage{enumitem,tabularx,fancyhdr,needspace}
\usepackage[unicode,hidelinks]{hyperref}
\hypersetup{pdfnewwindow=true,pdftitle={Econometrics 25117: Midterm 1 answer key with lecture references}}
\linespread{1}
\setlength{\parindent}{0pt}
\setlength{\parskip}{4pt}
\setlength{\abovedisplayskip}{5pt}
\setlength{\belowdisplayskip}{5pt}
\setlength{\emergencystretch}{1em}
\raggedbottom
\pagestyle{fancy}
\fancyhf{}
\renewcommand{\headrulewidth}{0pt}
\renewcommand{\footrulewidth}{0pt}
\fancyfoot[C]{\normalfont\thepage}
\pdfcatalog{/ViewerPreferences << /Duplex /DuplexFlipLongEdge >>}
\newcommand{\E}{\operatorname{E}}
\newcommand{\Var}{\operatorname{Var}}
\newcommand{\Cov}{\operatorname{Cov}}
\newcommand{\Corr}{\operatorname{Corr}}
\newcommand{\heading}[1]{\par\medskip\noindent #1\par\nobreak\smallskip}
% External PDF destinations use 1-based page values; hyperref writes a GoToR
% action with the corresponding 0-based PDF destination. Paths are relative
% to the final answer-key PDF in Exams/, not to the compilation directory.
\newcommand{\lecturepage}[2]{\href[page=#2]{../Lecture#1v2.pdf}{Lecture~#1, PDF p.~#2}}
\newcommand{\seminarpage}[2]{\href[page=#2]{../seminars/seminar_lecture#1.pdf}{Seminar~#1, PDF p.~#2}}
\newcommand{\lecturetext}[4]{\par\smallskip\noindent
  \lecturepage{#1}{#2}, ``#3'':\par\nobreak
  \noindent``#4''\par}
\newcommand{\seminartext}[4]{\par\smallskip\noindent
  Supplemental practice: \seminarpage{#1}{#2}, ``#3'':\par\nobreak
  \noindent``#4''\par}
\makeatletter
\renewcommand{\section}{\@startsection{section}{1}{\z@}%
  {6pt}{3pt}{\normalfont\fontsize{12}{14}\selectfont}}
\makeatother
\begin{document}
\begin{center}
{\fontsize{12}{14}\selectfont Econometrics (25117): Midterm 1, 2026--2027}\\
Worked answer key, explanations and lecture references: 150-minute midterm\\
\normalsize Monday, 9 November 2026
\end{center}

The maximum is 100 points; divide the raw total by 10 for the grade out of 10.
Award method credit for correct reasoning with an arithmetic error, and avoid
penalizing the same carried-forward error twice. Accept equivalent notation,
exact fractions and reasonable rounding. The solutions below use the
large-sample standard normal approximation, not an exact Student $t$ reference.

Each answer below includes the lecture sentences or formulas used to solve it.
Page numbers count PDF pages from 1, including the title page; they coincide
with the printed slide numbers in Lectures 1--3. The page references are
clickable links to the course PDFs. Keep this key in \textnormal{Exams/}, with the
lecture PDFs in the parent course folder and seminar PDFs in
\textnormal{seminars/} under that parent folder.

Quotation marks identify lecture wording. A sentence marked ``Application''
explains a calculation or deduction for this exam; it is not presented as a
quotation from the slides. Supplemental seminar references identify matching
practice where useful. The exam data differ from the examples in the lectures.

\section*{Multiple-choice questions}
\begin{center}
\begin{tabular}{c|cccccccc}
Question & 1 & 2 & 3 & 4 & 5 & 6 & 7 & 8 \\ \hline
Answer & b & c & b & a & d & b & d & c
\end{tabular}
\end{center}
Four points per correct option; zero for a wrong or blank response.
\begin{enumerate}[label=\arabic*.,leftmargin=*,itemsep=3pt,topsep=3pt]
\item Answer (b). $\E[Y_1]=\mu$, so the estimator is unbiased. Its distribution never
concentrates at $\mu$ as $n$ grows: the same nondegenerate random variable is
used for every $n$. In particular, for some $\varepsilon>0$,
$\Pr(|Y_1-\mu|>\varepsilon)>0$ remains unchanged. Thus it is inconsistent.
The positive variance ensures nondegeneracy. Option (a) incorrectly claims
consistency; (c) and (d) incorrectly claim bias, and (c) also incorrectly claims
consistency. Merely observing a nonvanishing variance is not a general proof
of inconsistency for arbitrary estimator sequences; here the unchanged
distribution supplies the argument.
% Lecture1v2.tex:601,706--707; PDF pages21,27.
\lecturetext{1}{21}{Random Sampling}{Then, the expected value of every draw is the population mean and the variance of every draw is the population variance.}
\lecturetext{1}{27}{Estimators}{$\hat\mu_Y$ is \textit{unbiased} if $\mathbb{E}(\hat\mu_Y)=\mu_Y$}
The same slide defines consistency by
$\hat\mu_Y\xrightarrow{p}\mu_Y$.
Application: substitute $Y_1$ into these definitions. Its expectation is right,
but its unchanged distribution does not concentrate at $\mu$ as $n$ grows.
This particular counterexample is an application of the definitions, rather
than a worked numerical example on that slide.
\item Answer (c). Independence implies $\E[XY]=\E[X]\E[Y]$ and hence zero
covariance. Options (a) and (b) reverse implications that do not hold in
general. For example, let $X$ be uniform on $[-1,1]$ and $Y=X^2$. Symmetry
gives $\Cov(X,Y)=0$, but $Y$ depends on $X$ and $\E[Y\mid X]=X^2$ is not
constant. Option (d) is false because perfect positive correlation between
nondegenerate variables means an exact increasing linear relationship, not
independence.
% Lecture1v2.tex:567--568; PDF page19. The quadratic graphic is on page20.
\lecturetext{1}{19}{Covariance}{When $X$ and $Y$ are independent, $\sigma_{XY}=0$. The converse is not true in general: zero covariance does not necessarily imply independence.}
The quadratic plot in panel (d) of \lecturepage{1}{20} illustrates a nonlinear
relationship with zero correlation. Application: the $Y=X^2$ example above
also shows why zero covariance need not make the conditional mean constant.
The uniform-distribution calculation is an added illustration of the slide's
point, not a quotation of a worked lecture example.
\item Answer (b). The CLT is
$\sqrt n(\bar Y-\mu)/\sigma\xrightarrow{d}N(0,1)$.
The factor $\sqrt n$ offsets the shrinking spread of the sample mean.
Option (a) incorrectly makes an exact finite-sample normality claim without
assuming a normal population. In (c), the expression without $\sqrt n$
converges in probability to zero. Option (d) is false because the law of
large numbers gives $\bar Y\xrightarrow{p}\mu$, not necessarily zero.
% Lecture1v2.tex:663,667,671--673,639,645--647; PDF pages25,23.
\lecturetext{1}{25}{Asymptotic properties ($n\rightarrow\infty$) of the sampling distribution (2/2)}{If $Y_1,Y_2,\ldots,Y_n$ are drawn randomly from the same population (i.i.d.) and large outliers are unlikely, then $\overline Y$ is approximately distributed as $\mathcal N(\mu_Y,\sigma_Y^2/n)$ as $n$ grows}
The CLT box on this page states
\[
\frac{\overline Y-\mu_Y}{\sigma_Y/\sqrt n}
\xrightarrow{d}\mathcal N(0,1).
\]
\lecturetext{1}{25}{Asymptotic properties ($n\rightarrow\infty$) of the sampling distribution (2/2)}{$\overline Y\sim\mathcal N(\mu_Y,\sigma_{\overline Y}^2)$ exactly when $Y\sim\mathcal N(\mu_Y,\sigma_Y^2)$, and approximately when $Y$ follows another distribution}
\lecturetext{1}{23}{Asymptotic properties ($n\rightarrow\infty$) of the sampling distribution (1/2)}{In this case, we say that $\overline Y$ converges in probability to $\mathbb E(Y)$}
Application: the exam does not assume a normal population, so the exact
finite-sample claim in (a) is not guaranteed. The CLT formula is algebraically
the same as (b). Convergence in probability concerns getting arbitrarily
close to the mean: $\Pr(|\bar Y-\mu|>\varepsilon)\to0$ for every
$\varepsilon>0$; it does not require exact equality.
\item Answer (a). A Type I error rejects a true null. Option (b) describes a
Type II error: failing to reject a false null. Options (c) and (d) describe
correct decisions, so neither is an error.
% Lecture2v2.tex:406,419; PDF page14.
\lecturetext{2}{14}{Type I and Type II Errors in Hypothesis Testing}{Occurs when we reject a null hypothesis that is actually true.}
This is the sentence in the Type I box. The Type II box on the same page says:
``Occurs when we fail to reject a null hypothesis that is actually false.''
Application: these are exactly the descriptions in options (a) and (b),
respectively; only the first answers the Type I question.
\item Answer (d). Coverage is a repeated-sampling property of the interval procedure.
In the frequentist framework the parameter is fixed, and the interval is random
before the sample is observed. With the approximation used here, coverage is
approximately 95\%. Option (a) assigns a probability to the fixed parameter
after observing the interval. Option (b) confuses uncertainty about a mean
with dispersion of individual observations. Option (c) confuses interval
coverage with exact equality between sample and population means.
% Lecture2v2.tex:546,557,562--568; PDF pages20--21.
\lecturetext{2}{20}{What are Confidence Intervals (CIs)?}{They provide a range of plausible values for a population parameter, not just a single point estimate.}
\lecturetext{2}{21}{Computing CIs}{For $Z\sim\mathcal N(0,1)$, choose $c$ so that $P(-c\leq Z\leq c)=1-\alpha$.}
The slide gives the interval $\overline y\pm c\times SE(\overline Y)$.
Application: with $c=1.96$, the large-sample probability statement can be
rewritten as
\[
\Pr\!\left(\mu_Y\in
[\overline Y-1.96SE(\overline Y),\,
 \overline Y+1.96SE(\overline Y)]\right)\approx0.95.
\]
The interval varies across samples while $\mu_Y$ stays fixed. This explains
the repeated-sampling wording in (d). The slides give the probability
construction; they do not print that exact repeated-sampling sentence.
\item Answer (b). The two residual identities are the sample OLS first-order
conditions when an intercept is included, regardless of the population
relationship between $X_i$ and $U_i$. Option (a) confuses these mechanical
sample identities with the population condition $\E[U_i\mid X_i]=0$.
Option (c) equates the unobserved population error with the fitted residual,
which generally differ because coefficients are estimated. Option (d) would
require a perfect fit; OLS makes residuals sum to zero, not each residual
equal zero.
% Lecture3v2.tex:270,274,276,330--331,605; Seminar3:407.
\lecturetext{3}{7}{The Ordinary Least Squares (OLS) Estimator}{The first-order conditions are}
\[
\sum_{i=1}^n(Y_i-\hat\beta_0-\hat\beta_1X_i)=0,
\qquad
\sum_{i=1}^n(Y_i-\hat\beta_0-\hat\beta_1X_i)X_i=0.
\]
Using $\hat U_i=Y_i-\hat Y_i$ from \lecturepage{3}{9} gives the two
residual identities in option (b).
\lecturetext{3}{19}{The conditions for causality}{The error term $U_i$ has conditional mean $0$ given $X_i$: $\mathbb E[U_i\mid X_i]=0$;}
Application: this is a population restriction, separate from the sample
equations used to fit the line.
\seminartext{3}{11}{Solution 3: Normal equations and squared errors (continued)}{These sample fitting identities do not establish causality.}
\item Answer (d). At $n=100$, the standard deviation of the sample mean is
$\sigma_Y/10$; at $n=400$, it is $\sigma_Y/20$. The ratio is $1/2$, so the
standard deviation is halved. Option (a) gives the change in the variance,
which is divided by four. Options (b) and (c) incorrectly say that uncertainty
increases or stays unchanged as the sample size quadruples.
% Lecture1v2.tex:625; PDF page22.
\lecturetext{1}{22}{Sample Distribution Mean and Dispersion}{The standard deviation of the sampling distribution is therefore given by $\frac{\sigma_Y}{\sqrt n}$}
Application: the ratio of the new standard deviation to the old one is
$\sqrt{100}/\sqrt{400}=10/20=1/2$. The population distribution has not changed;
the sample mean has become more precise.
\item Answer (c). $SER^2=SSR/(n-2)$ accounts for the two fitted coefficients,
while $RMSE^2=SSR/n$. Both describe the size of fitted residuals; neither is
the standard error of the slope. Option (a) reverses the denominators.
Option (b) is false because $SER>RMSE$ for positive $SSR$ and finite $n>2$.
Option (d) uses $n-1$, the denominator for a sample variance around an estimated
mean, rather than $n-2$ for this regression or $n$ for $RMSE$.
% Lecture3v2.tex:528,532,534,541; PDF page16.
\lecturetext{3}{16}{SER and RMSE: What Is the Difference?}{Both measures summarize the typical size of the OLS residuals.}
The two displayed definitions on that page are
\[
SER=\sqrt{\frac{SSR}{n-2}},\qquad RMSE=\sqrt{\frac{SSR}{n}}.
\]
The slide explains: ``The two degrees of freedom are used to estimate
$\beta_0$ and $\beta_1$.''
Application: these are the formulas in (c); for the same positive $SSR$,
dividing by the smaller denominator makes $SER$ larger than $RMSE$.
\end{enumerate}

\Needspace{7.5cm}
\section*{Long-answer questions}
\heading{Question 9: joint probabilities and dependence (10 points)}
\begin{enumerate}[label=(\alph*),leftmargin=*,itemsep=4pt,topsep=2pt]
\item $\Pr(D=1)=0.24+0.36=0.60$ and
$\Pr(Y=1)=0.04+0.36=0.40$. Each marginal probability sums the joint cells
over the other variable. Award 1 point each.
% Lecture1v2.tex:381--383; PDF page13.
The ``Marginal Probability distribution'' formula on
\lecturepage{1}{13}, ``Joint Probability Distribution $\Pr(X=x,Y=y)$'', is
\[
\Pr(Y=y)=\sum_i\Pr(X=x_i,Y=y).
\]
Application: sum down a column for the marginal of $D$ and across a row for
the marginal of $Y$. The result uses the exam's cells, not the lecture's
different joint-distribution example.
\item Because $Y$ is binary,
\[
\E[Y\mid D=0]=\frac{0.04}{0.40}=0.10,\qquad
\E[Y\mid D=1]=\frac{0.36}{0.60}=0.60.
\]
For a zero--one variable, its mean is the probability of being one; the
denominator restricts attention to the group being conditioned on.
The law of iterated expectations then gives
$\E[Y]=0.10(0.40)+0.60(0.60)=0.40$.
The overall mean is a weighted average of the two group means, with the
population group shares as weights.
Award 1 point for each conditional mean and 1 point for the weighted average.
% Lecture1v2.tex:385--387,442,455--463,489; PDF pages13--16.
The conditional-probability formula on \lecturepage{1}{13} is
$\Pr(Y=y\mid X=x)=\Pr(X=x,Y=y)/\Pr(X=x)$.
\lecturetext{1}{14}{Working with the joint distribution}{Same cell, different denominator: conditioning changes the reference group.}
\lecturetext{1}{16}{Calculating the law of iterated expectations}{Because $Y$ is binary, $\mathbb E(Y\mid X=x)=\Pr(Y=1\mid X=x)$.}
The ``Law of Iterated Expectations'' formula on \lecturepage{1}{15} is
\[
\mathbb E(Y)=\mathbb E[\mathbb E(Y\mid X)]
\equiv\sum_i\mathbb E(Y\mid X=x_i)\Pr(X=x_i).
\]
Application: replace $X$ by $D$, calculate the two within-group probabilities,
and weight them by the population shares $0.40$ and $0.60$.
\item The product $DY$ equals one only when both variables equal one. Thus
$\E[DY]=\Pr(D=1,Y=1)=0.36$, so
\[
\Cov(D,Y)=0.36-(0.60)(0.40)=0.12.
\]
For a binary variable with probability $p$ of being one, its variance is
$p(1-p)$. Both variances therefore equal $0.60(0.40)=0.24$. Consequently,
\[
\Corr(D,Y)=\frac{0.12}{\sqrt{0.24\cdot0.24}}=0.50.
\]
Correlation divides covariance by the product of the two standard deviations,
removing the units. The positive sign means that workshop attendance and
employment are positively associated; it does not by itself establish a
causal effect.
Award 1.5 points for the covariance and 1.5 for the correlation, including
the variance calculation.
% Lecture1v2.tex:560--564,577--578,294--300; PDF pages19,20,10.
\lecturetext{1}{19}{Covariance}{The covariance between $X$ and $Y$, $\sigma_{XY}$ or $\operatorname{Cov}(X,Y)$, is the average product of how two variables differ from their means. It measures how $X$ and $Y$ move together.}
\lecturetext{1}{20}{Correlation}{To clarify the relationship, covariance is normalized by the variables' standard deviations.}
The correlation formula on that slide is
$\rho_{XY}=\sigma_{XY}/(\sigma_X\sigma_Y)\in[-1,1]$.
The moment definitions on \lecturepage{1}{10}, ``Statistical Moments'', include
$\operatorname{Var}(Y)=\mathbb E[(Y-\mu_Y)^2]$.
Application: expanding the covariance definition gives
$\Cov(D,Y)=\E[DY]-\E[D]\E[Y]$.
For a zero--one variable $B$, $B^2=B$, so its variance is
$\E[B^2]-\E[B]^2=p-p^2=p(1-p)$. These algebraic steps give the covariance
and both variances used above.
\item They are not independent: $\Pr(D=1,Y=1)=0.36$ differs from
$\Pr(D=1)\Pr(Y=1)=0.24$. Equivalently,
$\Pr(Y=1\mid D=1)=0.60\ne0.40=\Pr(Y=1)$.
Award 1 point for the conclusion and 1 point for a valid probability comparison.
% Lecture1v2.tex:521--526; PDF page17.
\lecturetext{1}{17}{Independence}{The joint probability distribution of $X$ and $Y$ can be expressed as the product of their individual probability distributions}
The same slide writes the independence criterion as
\[
\Pr(Y=y\mid X=x)=\Pr(Y=y)
\quad\Longleftrightarrow\quad
\Pr(X=x,Y=y)=\Pr(Y=y)\Pr(X=x).
\]
Application: one cell violating the product equality is enough to rule out
independence. Here $0.36\ne0.24$.
\end{enumerate}

\heading{Question 10: comparing two means (11 points)}
\begin{enumerate}[label=(\alph*),leftmargin=*,itemsep=4pt,topsep=2pt]
\item The estimated difference is $\widehat\Delta=74-68=6$ points. Under the
stated independence,
\[
SE(\widehat\Delta)=\sqrt{\frac{12^2}{100}+\frac{16^2}{100}}
=\sqrt{1.44+2.56}=2\text{ points}.
\]
Each sample mean has estimated variance $s_g^2/n_g$. Independence makes the
variance of their difference the sum of these variances; the standard error
is the square root of that sum.
Award 1 point for the estimate and 1 point for the standard error.
% Lecture2v2.tex:594--603; PDF page22.
\lecturetext{2}{22}{Mean Difference Testing: Comparing Groups 1/2}{Assume two independent samples of sizes $n_a$ and $n_b$, with i.i.d. observations within each sample. Thus $\operatorname{Cov}(\overline Y^a,\overline Y^b)=0$.}
The slide says: ``Estimate the standard error using the sample standard
deviations, $s_Y^a$ and $s_Y^b$:'' and gives
\[
SE(\overline Y^a-\overline Y^b-d_0)
=\sqrt{\frac{(s_Y^a)^2}{n_a}+\frac{(s_Y^b)^2}{n_b}}.
\]
Application: subtracting the fixed null value $d_0$ does not change the
variance. Insert the two sample SDs and sizes to obtain $SE=2$.
\item Test $H_0:\Delta=2$ against $H_1:\Delta\ne2$.
The observed test statistic is
$t_{\mathrm{obs}}=(6-2)/2=2$.
Award 0.5 points for each hypothesis and 2 points for the correctly centered
and standardized statistic. Using $6/2=3$ tests the wrong null.
% Lecture2v2.tex:587,612--614; PDF pages22--23.
\lecturetext{2}{23}{Mean Difference Testing: Comparing Groups 2/2}{Once we have $SE(\overline Y^a-\overline Y^b)$, we can easily compute the corresponding t-statistic}
\[
t=\frac{(\overline y^a-\overline y^b)-d_0}
{SE(\overline Y^a-\overline Y^b)}.
\]
The hypotheses on \lecturepage{2}{22} allow a general null difference $d_0$.
Application: set $d_0=2$ here. The numerator is the distance between the
estimated difference 6 and the claimed difference 2.
\item For the two-sided test, compare the absolute statistic with each
critical value. At 5\%, $|2|>1.96$, so reject $H_0$. At 1\%,
$|2|<2.576$, so do not reject $H_0$. A stricter significance level requires
stronger evidence against the null. Award 1 point for each comparison
and its correct decision. No probability-table lookup is required.
% Lecture2v2.tex:616,562--568; Seminar2:331,335--338.
\lecturetext{2}{23}{Mean Difference Testing: Comparing Groups 2/2}{We can reject that $\mu_Y^a-\mu_Y^b=d_0$ at the 5\%-level when $\mid t\mid>1.96$}
\seminartext{2}{10}{Solution 2: Decisions and the meaning of .0164 (continued)}{Reject $H_0$ when the p-value is at most $\alpha$, or equivalently when $|t^{\mathrm{obs}}|$ reaches the two-sided critical value.}
The seminar table prints critical values 1.960 at 5\% and 2.576 at 1\%.
The table on \lecturepage{2}{21} prints the 1\% value rounded to 2.58.
Application: use the more precise value supplied in the exam; either rounding
gives the same decision for $|t|=2$.
\item The 95\% interval is
\[
6\pm1.96(2)=[2.08,9.92]\text{ points}.
\]
It excludes the null value 2, consistently with rejection by the two-sided
5\% test using the same normal approximation. Award 0.5 points for the
formula, 0.5 for each endpoint, 0.25 for checking the value 2, and 0.25 for
the test--interval connection. The interval concerns the population difference,
not the spread of individual scores.
% Lecture2v2.tex:617--621; PDF page23.
\lecturetext{2}{23}{Mean Difference Testing: Comparing Groups 2/2}{Similarly, we can construct the confidence interval for $\mu_Y^a-\mu_Y^b$}
\[
CI_{\mu_Y^a-\mu_Y^b}
=(\overline y^a-\overline y^b)\pm1.96\times
SE(\overline Y^a-\overline Y^b).
\]
The same slide states: ``We can reject that $\mu_Y^a-\mu_Y^b=d_0$ at the
5\%-level when $d_0\notin CI_{\mu_Y^a-\mu_Y^b}$.''
Application: the excluded value is 2, which is the null being tested;
checking only whether zero belongs to the interval would answer a different
test question.
\item With $m$ students in each independent group, the planned standard
error is $\sqrt{12^2/m+16^2/m}=20/\sqrt m$. The planned 95\% interval
half-width is therefore
\[
h(m)=1.96\frac{20}{\sqrt m}=\frac{39.2}{\sqrt m}.
\]
Requiring $h(m)\leq2$ gives $\sqrt m\geq19.6$, or
$m\geq19.6^2=384.16$. Sample size must be an integer and must meet the
inequality, so round up: $m_{\min}=385$ students per group, or 770 in total.
For comparison, $m=384$ gives a half-width just above 2. These are planning
calculations using the supplied standard deviations; the realized interval
width can differ when standard deviations are estimated from the new data.
Award 1 point for the half-width inequality and bound, 0.5 for rounding up to
385 per group, and 0.5 for the total of 770.
% Lecture2v2.tex:594--603,617--619; Seminar2:538--547,561--562.
Application of the independent-group variance formula on
\lecturepage{2}{22} and the interval formula on \lecturepage{2}{23}:
replace both group sizes by $m$ and solve the half-width inequality.
This two-group planning calculation combines the lecture formulas; it is
not a verbatim worked example in the lecture.
\seminartext{2}{20}{Solution 6: Sample size for a target half-width}{For $m_n\leq m$, solve and round up:}
That seminar uses $m$ for a target half-width and $n$ for the sample size,
and gives $n_{\min}=\lceil(1.96\sigma/m)^2\rceil$ for one mean. In this exam,
$m$ instead denotes students per group. Apply the same solve-and-round-up
method to the two-group standard error derived above.
The seminar also states: ``Realized precision depends on how well the
planning value describes the population.''
\end{enumerate}

\Needspace{4\baselineskip}
\heading{Question 11: simple OLS (11 points)}
\begin{enumerate}[label=(\alph*),leftmargin=*,itemsep=4pt,topsep=2pt]
\item For simple OLS with an intercept, divide the centered cross-product
sum by the centered sum of squares of $X$:
\[
\widehat\beta_1=
\frac{\sum_i(x_i-\bar x)(y_i-\bar y)}{\sum_i(x_i-\bar x)^2}
=\frac{50}{100}=0.5,\qquad
\widehat\beta_0=70-0.5(4)=68.
\]
The intercept follows from the fact that the fitted line passes through
$(\bar x,\bar y)$.
The fitted line is $\widehat y_i=68+0.5x_i$: one additional weekly study hour
is associated with 0.5 additional predicted score points. This is not yet a
causal statement. Award 0.5 points each for the slope, intercept, fitted line
and interpretation.
% Lecture3v2.tex:307--316,184; PDF pages8,3.
\lecturetext{3}{8}{The Ordinary Least Squares (OLS) Estimator}{The intercept adjusts the average outcome for the fitted slope.}
The displayed coefficient formulas are
\[
\hat\beta_0=\overline Y-\hat\beta_1\overline X,\qquad
\hat\beta_1=\frac{\sum_i(Y_i-\overline Y)(X_i-\overline X)}
{\sum_i(X_i-\overline X)^2}=\frac{s_{XY}}{s_X^2}.
\]
The same slide says: ``The slope compares the comovement of $X$ and $Y$
with the variation in $X$.''
\lecturetext{3}{3}{The regression line}{The slope is the difference in the expected values of Y, for two values of X that differ by one unit}
Application: use the supplied centered sums in the estimator formula and
read the fitted slope as the change in predicted score per additional hour.
The estimated-line interpretation alone makes no causal claim.
\item $\widehat y_i=68+0.5(6)=71$ and
$\widehat u_i=72-71=1$. The fitted score is the point on the estimated line
at six hours. The residual is observed minus fitted, so this student scores
one point above the prediction. Award 0.5 points for the fitted score and
0.5 for the residual.
% Lecture3v2.tex:330--331; Seminar3:379.
The formulas on \lecturepage{3}{9}, ``Back to the Linear Regression Model'', are
\[
\hat Y_i=\hat\beta_0+\hat\beta_1X_i,\qquad
\hat U_i=Y_i-\hat Y_i.
\]
\seminartext{3}{10}{Solution 3: Fitted scores and vertical residuals}{Residuals use the estimated sample line. Population errors use the true population coefficients.}
Application: substitute six hours into the fitted line, then subtract its
prediction from the observed score. Reversing the subtraction would give
the wrong residual sign.
\item Since OLS includes an intercept,
\[
ESS=\widehat\beta_1^2\sum_i(x_i-\bar x)^2=0.5^2(100)=25,
\qquad SSR=TSS_Y-ESS=425-25=400.
\]
Thus $R^2=25/425=1/17\approx0.0588$ and
$SER=\sqrt{400/(102-2)}=2$ score points.
The fitted values account for about 5.88\% of the sample variation in scores.
$SER$ summarizes residual dispersion in score points, using $n-2$ because
the fitted line estimates an intercept and a slope. It is not a coefficient
standard error. Award 1 point each for $ESS$, $SSR$, $R^2$ and $SER$.
% Lecture3v2.tex:434,421--422,528,532,547; Seminar3:597.
\lecturetext{3}{13}{$R^2$}{The $R^2$ is the ratio of the sample variance of $\hat Y$ to the sample variance of $Y$}
This page writes $R^2$ as the explained sum of squares divided by the total
sum of squares, and also as one minus the residual sum of squares divided
by the total sum of squares.
\lecturetext{3}{12}{Goodness of Fit}{The standard error of the regression ($SER$) measures the magnitude of a typical regression residual in the units of Y.}
The definition $SER=\sqrt{SSR/(n-2)}$ appears on \lecturepage{3}{16}.
\seminartext{3}{19}{Solution 6: Decomposing the observed score variation}{Thus $TSS=ESS+SSR$: $32=25.6+6.4$.}
Application: the seminar numbers illustrate the decomposition. For this exam,
$\hat y_i-\bar y=0.5(x_i-\bar x)$, so $ESS=0.5^2(100)=25$ and the
corresponding decomposition is $425=25+400$.
\item Sample correlation can be calculated directly from centered sums:
\[
r_{XY}=\frac{\sum_i(x_i-\bar x)(y_i-\bar y)}
{\sqrt{\sum_i(x_i-\bar x)^2\sum_i(y_i-\bar y)^2}}
=\frac{50}{\sqrt{100\cdot425}}
=\frac{1}{\sqrt{17}}\approx0.2425.
\]
The sample covariance and both sample variances would each use the same
$n-1$ divisor; it cancels in this correlation formula. Squaring gives
$r_{XY}^2=1/17\approx0.0588=R^2$. This equality holds for simple OLS with
an intercept. Correlation retains the positive sign of the association,
whereas $R^2$ records its squared strength. Award 1 point for the correlation
and 1 point for verifying the equality numerically.
% Lecture2v2.tex:748,750; Lecture3v2.tex:309--312,436--437.
\lecturetext{2}{28}{The California Test Score Data Set}{The sample correlation, $r_{XY}=\frac{s_{XY}}{s_Xs_Y}$}
The same page gives
$s_{XY}=(n-1)^{-1}\sum_i(X_i-\overline X)(Y_i-\overline Y)$.
Application: the $n-1$ terms cancel when this covariance is divided by the
two sample standard deviations, giving the centered-sum correlation above.

The identity $r^2=R^2$ is an algebraic consequence of the OLS formulas on
\lecturepage{3}{8} and the fit definition on \lecturepage{3}{13}; it is not
printed there as a separate sentence. Writing $S_{XX},S_{XY},S_{YY}$ for the
three centered sums supplied in the question gives
\[
ESS=\hat\beta_1^2S_{XX}=\frac{S_{XY}^2}{S_{XX}},\qquad
R^2=\frac{S_{XY}^2}{S_{XX}S_{YY}}=r_{XY}^2.
\]
This derivation requires simple OLS with an intercept, as in the exam.
\item Substitute $X_i=M_i/60$ into the original fitted line:
\[
\widehat y_i=68+0.5\left(\frac{M_i}{60}\right)
=68+\frac{1}{120}M_i.
\]
The intercept remains 68, and the slope is $1/120\approx0.00833$ score
points per additional study minute. A minute is $1/60$ of an hour, so its
slope is $1/60$ of the hourly slope. The conversion changes the numerical
regressor units, but it does not change any fitted score or residual.
Consequently $ESS$, $SSR$ and $TSS_Y$ remain the same, giving
$R^2=1/17\approx0.0588$ and $SER=2$ score points again. Award 0.5 points
each for the intercept, slope, unchanged $R^2$ and unchanged $SER$.
% Lecture3v2.tex:330--331,421,547; Seminar3:451.
Application of the fitted-value formula on \lecturepage{3}{9}: expressing
the same study time in minutes means $X_i=M_i/60$. Substitute that identity
without changing the fitted score.
\seminartext{3}{13}{Solution 4: Shares, percentage points, and predicted scores}{Changing the units changes the coefficient's numerical scale and preserves the predictions.}
The seminar example converts a share to percentage points; the exam applies
the same substitution method to hours and minutes.
\lecturetext{3}{16}{SER and RMSE: What Is the Difference?}{Both are measured in the units of $Y$. They describe the typical size of the part of $Y$ that the regression does not explain.}
Application: the outcome units and every residual are unchanged. Therefore
the residual sum of squares and $SER$ are unchanged. The explained and total
sums of squares are also unchanged, so $R^2$ stays the same, consistently
with its unitless interpretation on \lecturepage{3}{12}.
\end{enumerate}


\heading{Additional paper-and-pencil problems: Questions 12--14}

\heading{Question 12: Lesson 1---a small library population (12 points)}
\begin{enumerate}[label=(\alph*),leftmargin=*,itemsep=4pt,topsep=2pt]
\item There is one student with $Y=0$, six with $Y=2$ and one with $Y=4$.
Each student has probability $1/8$ of being selected, so
\[
\Pr(Y=0)=\frac18,\qquad \Pr(Y=2)=\frac68=\frac34,\qquad
\Pr(Y=4)=\frac18.
\]
The probabilities add to one. Accumulating the probabilities from left to
right gives the complete c.d.f.:
\[
F_Y(y)=
\begin{cases}
0,&y<0,\\
1/8,&0\leq y<2,\\
7/8,&2\leq y<4,\\
1,&y\geq4.
\end{cases}
\]
Award 1.5 points for the marginal distribution and 1.5 points for the complete c.d.f.,
including its values at the jumps.
\lecturetext{1}{7}{Some Revision}{The cumulative distribution function (c.d.f.) is the probability that the random variable is less than or equal to a particular value}
The marginal-distribution formula on \lecturepage{1}{13},
``Joint Probability Distribution $\Pr(X=x,Y=y)$'', is
$\Pr(Y=y)=\sum_i\Pr(X=x_i,Y=y)$.
Application: add the two programme counts in each row and divide by eight,
then sum the probabilities at all outcomes no larger than $y$.
\seminartext{1}{5}{Solution 1: The earnings c.d.f. (continued)}{At each jump, the filled dot gives the value of the right-continuous c.d.f.}
Application: use the same step-function construction, with this problem's
jumps at 0, 2 and 4. In particular, $F_Y(0)=1/8$, $F_Y(2)=7/8$ and $F_Y(4)=1$.
\item Weight each value and its square by its probability:
\[
\E[Y]=0\!\left(\frac18\right)+2\!\left(\frac34\right)
       +4\!\left(\frac18\right)=\frac32+\frac12=2,
\]
\[
\E[Y^2]=0^2\!\left(\frac18\right)+2^2\!\left(\frac34\right)
       +4^2\!\left(\frac18\right)=3+2=5.
\]
Therefore
\[
\Var(Y)=\E[Y^2]-(\E[Y])^2=5-2^2=1.
\]
A direct check gives
$\Var(Y)=(0-2)^2/8+(2-2)^2(3/4)+(4-2)^2/8=1$.
The mean is two books and the variance is one book squared. These are moments
of the specified eight-student population, not estimates from a sample;
there is no $n-1$ correction. Award 1 point for each requested moment.
\lecturetext{1}{10}{Statistical Moments}{The mean (or expected value) indicates the long-run average of a random variable $Y$:}
That slide gives $\mathbb E(Y)=\sum_i y_i\Pr(Y=y_i)$ and
$\operatorname{Var}(Y)=\mathbb E[(Y-\mu_Y)^2]$.
The variance shortcut on \lecturepage{1}{11}, ``Calculating variance two ways'',
is $\operatorname{Var}(Y)=\mathbb E(Y^2)-[\mathbb E(Y)]^2$.
Application: substitute the marginal probabilities from part (a).
\item Each programme contains four students, so
$\Pr(G=\mathrm A)=\Pr(G=\mathrm B)=4/8=1/2$.
Programme A contains one student borrowing zero books, two borrowing two and
one borrowing four. Programme B has four students, all borrowing two. Thus
\[
\E[Y\mid G=\mathrm A]=\frac{0(1)+2(2)+4(1)}4=\frac84=2,
\qquad
\E[Y\mid G=\mathrm B]=\frac{2(4)}4=2.
\]
Iterated expectations weights the two conditional means by their population
programme shares:
\[
\E[Y]=\E[Y\mid G=\mathrm A]\Pr(G=\mathrm A)
     +\E[Y\mid G=\mathrm B]\Pr(G=\mathrm B)
     =2\!\left(\frac12\right)+2\!\left(\frac12\right)=2.
\]
This agrees with part (b). Award 1 point for each conditional mean and
1 point for the correctly weighted recovery of the marginal mean.
The formulas on \lecturepage{1}{15}, ``Conditional and Iterated Expectations'',
are
\[
\mathbb E(Y\mid X=x)=\sum_i y_i\Pr(Y=y_i\mid X=x),\qquad
\mathbb E(Y)=\sum_i\mathbb E(Y\mid X=x_i)\Pr(X=x_i).
\]
\lecturetext{1}{14}{Working with the joint distribution}{Same cell, different denominator: conditioning changes the reference group.}
Application: within a programme, divide its counts by four rather than eight.
For the marginal mean, average the two programme means using their shares.
\item The variables are not independent. For example,
\[
\Pr(Y=2\mid G=\mathrm B)=\frac44=1,
\qquad \Pr(Y=2)=\frac68=\frac34.
\]
Knowing that the selected student is in programme B changes the probability
of borrowing two books. The conditional means happen to be equal, but the
conditional distributions differ: programme A includes zero and four books,
whereas programme B contains only two. Equality of conditional means alone
therefore does not imply independence.
Award 1.5 points for the conditional-versus-marginal probability comparison
and 1.5 points for the conclusion about independence and equal means.
The independence criterion on \lecturepage{1}{17}, ``Independence'', includes
$\Pr(Y=y\mid X=x)=\Pr(Y=y)$.
Application: the displayed failure at $y=2$ rules out independence. The
conclusion about equal means is a deduction from this problem's calculations,
not a quotation of a worked example in the lecture.
\end{enumerate}

\heading{Question 13: Lesson 2---from sample summaries to inference (12 points)}
\begin{enumerate}[label=(\alph*),leftmargin=*,itemsep=4pt,topsep=2pt]
\item The sample mean is
\[
\bar y=\frac{5200}{100}=52\text{ points}.
\]
The supplied raw-square identity gives
\[
S_{YY}=280300-100(52)^2=280300-270400=9900\text{ points}^2.
\]
Thus
\[
s_Y^2=\frac{9900}{100-1}=100\text{ points}^2,\qquad
s_Y=10\text{ points},\qquad
SE(\bar Y)=\frac{10}{\sqrt{100}}=1\text{ point}.
\]
The sample SD describes dispersion of individual scores; the SE describes
sampling uncertainty in the estimated mean. Award 0.5 points for the mean, 0.75 for the centered sum of squares,
0.75 for the sample variance, 0.5 for the sample SD, and 0.5 for the SE,
with appropriate units (3 points in total).
\lecturetext{2}{8}{Sample Variance and Sample Standard Deviation}{The sample variance, $s_Y^2$ is an estimator of the population variance $\sigma_Y^2$}
The page defines $s_Y^2=(n-1)^{-1}\sum_i(Y_i-\overline Y)^2$ and
$s_Y=\sqrt{s_Y^2}$. Application: the raw-square identity in the question
follows by expanding $(y_i-\bar y)^2$ and using $\sum_i y_i=n\bar y$;
it is an algebraic shortcut, not a sentence quoted from the lecture.
\lecturetext{2}{10}{Step 2: Estimating the Standard Error}{Replacing $\sigma_Y$ by $s_Y$ gives the estimated standard deviation of $\overline Y$:}
The displayed formula is $SE(\overline Y)=s_Y/\sqrt n$.
\lecturetext{2}{10}{Step 2: Estimating the Standard Error}{The standard error measures the sampling uncertainty of the estimator $\overline Y$.}
\item The hypotheses are $H_0:\mu_Y=50$ and $H_1:\mu_Y\ne50$.
The observed statistic is
\[
t_{\mathrm{obs}}=\frac{52-50}{1}=2.
\]
At 5\%, $|2|>1.96$, so reject $H_0$. At 1\%, $|2|<2.576$, so do not
reject $H_0$. This uses the stated large-sample normal approximation;
no exact Student $t$ reference distribution is being assumed.
Award 0.5 points for each hypothesis, 1 point for the correctly centered
and standardized statistic, and 0.5 for each correct critical-value
comparison and decision (3 points in total).
\lecturetext{2}{7}{P-Value for a One-Sample Mean Test}{We test a population mean against a hypothesized value:}
The hypotheses on that page are $H_0:\mu_Y=\mu_Y^0$ versus
$H_1:\mu_Y\ne\mu_Y^0$. Set $\mu_Y^0=50$ here.
The statistic displayed on \lecturepage{2}{13}, ``The t-statistic and its
use for hypothesis testing'', is
$t=(\overline y-\mu_Y^0)/SE(\overline Y)$.
\seminartext{2}{10}{Solution 2: Decisions and the meaning of .0164 (continued)}{Reject $H_0$ when the p-value is at most $\alpha$, or equivalently when $|t^{\mathrm{obs}}|$ reaches the two-sided critical value.}
The seminar table gives 1.960 at 5\% and 2.576 at 1\%. The table on
\lecturepage{2}{21} rounds the 1\% value to 2.58; that rounding gives the
same decision in this problem.
\item The 95\% normal approximate interval is
\[
52\pm1.96(1)=[50.04,53.96]\text{ points}.
\]
It excludes 50, in agreement with rejection by the two-sided 5\% test.
Do not round the lower endpoint to 50 before checking inclusion.
Award 0.75 points for the interval formula, 0.75 for each endpoint, and
0.75 for checking 50 and connecting the result to the test (3 points in total).
\lecturetext{2}{21}{Computing CIs}{Use $\overline y\pm c\times SE(\overline Y)$, adding and subtracting $c$ standard errors.}
The same slide says: ``Excluding 175 cm is equivalent to rejecting
$H_0:\mu_Y=175$ cm at 5\% (two-sided).'' Application: replace the lecture's
null value 175 by this problem's null value 50; it is excluded here.
\item Using the planning SD, the planned half-width is
\[
h(n)=1.96\frac{10}{\sqrt n}=\frac{19.6}{\sqrt n}.
\]
The required inequality is
\[
\frac{19.6}{\sqrt n}\leq0.98
\quad\Longleftrightarrow\quad
\sqrt n\geq20
\quad\Longleftrightarrow\quad n\geq400.
\]
The minimum integer is therefore 400 students. It is four times the original
sample size. With the SD unchanged, quadrupling $n$ halves the planned SE
from 1 to 0.5 and the planned half-width from 1.96 to 0.98. The actual
interval width can differ if the new survey's estimated SD differs from
the planning value. Award 1 point for the half-width inequality, 1 point for
the minimum sample size, and 1 point for the sample-size/precision comparison
(3 points in total).
\lecturetext{2}{9}{Step 1: The Standard Deviation of the Sample Mean}{This is smaller when $n$ is larger: averaging more observations makes the sample mean more precise.}
That page gives the theoretical SD $\sigma_Y/\sqrt n$. Combine it with
the interval construction on \lecturepage{2}{21} to form the half-width.
This planning inequality is an application of those lecture formulas.
\seminartext{2}{20}{Solution 6: Sample size for a target half-width}{For $m_n\leq m$, solve and round up:}
The seminar's formula is
$n_{\min}=\lceil(1.96\sigma/m)^2\rceil$, where $m$ denotes target
half-width. Insert $\sigma=10$ and $m=0.98$ to obtain exactly 400.
\end{enumerate}

\Needspace{7cm}
\heading{Question 14: Lesson 3---simple OLS from five observations (12 points)}
\begin{enumerate}[label=(\alph*),leftmargin=*,itemsep=4pt,topsep=2pt]
\item The means are $\bar x=15/5=3$ and $\bar y=35/5=7$.
Set $d_{xi}=x_i-3$ and $d_{yi}=y_i-7$. The centered calculations are:
\begin{center}
\begin{tabular}{c|rrrrrr}
$i$ & $x_i$ & $y_i$ & $d_{xi}$ & $d_{yi}$ & $d_{xi}^2$ & $d_{xi}d_{yi}$ \\ \hline
1 & 1 & 4 & $-2$ & $-3$ & 4 & 6 \\
2 & 2 & 4 & $-1$ & $-3$ & 1 & 3 \\
3 & 3 & 7 & 0 & 0 & 0 & 0 \\
4 & 4 & 8 & 1 & 1 & 1 & 1 \\
5 & 5 & 12 & 2 & 5 & 4 & 10 \\ \hline
Sum & 15 & 35 & 0 & 0 & 10 & 20
\end{tabular}
\end{center}
Therefore $S_{XX}=10$, $S_{XY}=20$,
$\widehat\beta_1=20/10=2$ and
$\widehat\beta_0=7-2(3)=1$.
The fitted line is $\widehat y_i=1+2x_i$. One additional weekly study hour is
associated with two additional predicted quiz points. This interpretation
of the fitted association alone does not establish a causal effect.
Award 0.5 points for each mean, 0.5 for each centered sum, 0.5 for each
coefficient, 0.5 for the fitted line, and 0.5 for the slope interpretation.
% Lecture3v2.tex:307--316; PDF page8.
\lecturetext{3}{8}{The Ordinary Least Squares (OLS) Estimator}{The intercept adjusts the average outcome for the fitted slope.}
The same page gives
\[
\hat\beta_1=\frac{\sum_i(Y_i-\overline Y)(X_i-\overline X)}
{\sum_i(X_i-\overline X)^2},\qquad
\hat\beta_0=\overline Y-\hat\beta_1\overline X.
\]
Its other bullet says: ``The slope compares the comovement of $X$ and $Y$
with the variation in $X$.'' Application: the centered table supplies the
numerator and denominator directly; the slope is measured in quiz points
per weekly study hour.
\Needspace{5.5cm}
\item Use $\widehat y_i=1+2x_i$ and
$\widehat u_i=y_i-\widehat y_i$:
\begin{center}
\begin{tabular}{c|rrrrr}
$i$ & $x_i$ & $y_i$ & $\widehat y_i$ & $\widehat u_i$ & $x_i\widehat u_i$ \\ \hline
1 & 1 & 4 & 3 & 1 & 1 \\
2 & 2 & 4 & 5 & $-1$ & $-2$ \\
3 & 3 & 7 & 7 & 0 & 0 \\
4 & 4 & 8 & 9 & $-1$ & $-4$ \\
5 & 5 & 12 & 11 & 1 & 5 \\ \hline
Sum & 15 & 35 & 35 & 0 & 0
\end{tabular}
\end{center}
Thus $\sum_i\widehat u_i=1-1+0-1+1=0$ and
$\sum_i x_i\widehat u_i=1-2+0-4+5=0$.
These are the sample OLS fitting identities for an intercept and one
regressor. Individual residuals can be positive, negative or zero even
though both sums equal zero.
Award 1 point for the five fitted scores, 1 point for the residuals, and
0.5 for each verified zero sum.
% Lecture3v2.tex:330--331,270,274,276; PDF pages9,7.
The formulas on \lecturepage{3}{9}, ``Back to the Linear Regression Model'', are
$\hat Y_i=\hat\beta_0+\hat\beta_1X_i$ and
$\hat U_i=Y_i-\hat Y_i$.
\lecturetext{3}{7}{The Ordinary Least Squares (OLS) Estimator}{The first-order conditions are}
\[
\sum_i(Y_i-\hat\beta_0-\hat\beta_1X_i)=0,\qquad
\sum_i(Y_i-\hat\beta_0-\hat\beta_1X_i)X_i=0.
\]
Application: use the estimated coefficients in the first two formulas and
add the residuals and the hour-weighted residuals to check the latter two.
\seminartext{3}{11}{Solution 3: Normal equations and squared errors (continued)}{These sample fitting identities do not establish causality.}
\Needspace{5.5cm}
\item The squared deviations are:
\begin{center}
\begin{tabular}{c|rrr}
$i$ & $(y_i-7)^2$ & $(\widehat y_i-7)^2$ & $\widehat u_i^2$ \\ \hline
1 & 9 & 16 & 1 \\
2 & 9 & 4 & 1 \\
3 & 0 & 0 & 0 \\
4 & 1 & 4 & 1 \\
5 & 25 & 16 & 1 \\ \hline
Sum & $TSS=44$ & $ESS=40$ & $SSR=4$
\end{tabular}
\end{center}
The decomposition is $44=40+4$. Hence
\[
R^2=\frac{ESS}{TSS}=\frac{40}{44}=\frac{10}{11}
=1-\frac4{44},\qquad
SER=\sqrt{\frac{SSR}{n-2}}=\sqrt{\frac4{5-2}}
=\frac2{\sqrt3}.
\]
The fitted line explains $10/11$ of the sample variation in quiz scores;
$1/11$ remains in the residuals. $SER$ is residual dispersion measured in
quiz points, with two fitted coefficients accounted for by $n-2$.
It is not a standard error of a coefficient. Accept
$\sqrt{4/3}$ or $2\sqrt3/3$ for the same exact $SER$.
Award 0.5 points each for $TSS$, $ESS$ and $SSR$, 0.75 for $R^2$ and
0.75 for $SER$. The displayed decomposition verifies the three sums;
carry forward one consistent arithmetic error without penalizing it twice.
% Lecture3v2.tex:434,436--437,532,534; PDF pages13,16.
\lecturetext{3}{13}{$R^2$}{The $R^2$ is the ratio of the sample variance of $\hat Y$ to the sample variance of $Y$}
That page displays $R^2=ESS/TSS=1-SSR/TSS$ with the corresponding squared
sum definitions.
\lecturetext{3}{16}{SER and RMSE: What Is the Difference?}{The two degrees of freedom are used to estimate $\beta_0$ and $\beta_1$.}
The $SER$ formula on that page is $SER=\sqrt{SSR/(n-2)}$.
\seminartext{3}{19}{Solution 6: Decomposing the observed score variation}{Thus $TSS=ESS+SSR$: $32=25.6+6.4$.}
Application: the seminar demonstrates the same identity with different
numbers. The exam calculations give $44=40+4$ and the exact residual
spread $2/\sqrt3$; neither a decimal approximation nor a probability table
is needed.
\item Because $X_i=M_i/60$, substitute into the fitted line:
\[
\widehat y_i=1+2\left(\frac{M_i}{60}\right)
=1+\frac1{30}M_i.
\]
The intercept is still 1 quiz point. The slope is $1/30$ of a quiz point per
additional study minute: 60 minutes correspond to the original two-point
change per study hour. This is a change of regressor units, so every fitted
score and residual stays the same. The outcome units, sample size and sums
of squares also stay the same. Therefore $R^2=10/11$ and
$SER=2/\sqrt3$ quiz points remain unchanged.
Award 0.5 points each for the intercept, minute slope, unchanged $R^2$ and
unchanged $SER$, with the invariance explanation included in the last two items.
% Seminar3:443,445,451; Lecture3v2.tex:330--331,436--437,532,547.
\seminartext{3}{13}{Solution 4: Shares, percentage points, and predicted scores}{Changing the units changes the coefficient's numerical scale and preserves the predictions.}
The seminar converts a share to percentage points. Application: here the
same substitution method converts hours to minutes; $X_i=M_i/60$ is the
reason the slope is divided by 60.
\lecturetext{3}{16}{SER and RMSE: What Is the Difference?}{Both are measured in the units of $Y$. They describe the typical size of the part of $Y$ that the regression does not explain.}
Application: the fitted scores, residuals and outcome units have not
changed. Thus the sums in the $R^2$ definition on \lecturepage{3}{13} and
the $SER$ definition on \lecturepage{3}{16} remain the same. The numerical
invariance after this specific hours-to-minutes conversion is a deduction,
not a separate verbatim lecture example.
\end{enumerate}

%\heading{Coverage and time budget (instructor only)}
%Questions 1--8: 12 minutes; Q9: 13; Q10: 20; Q11: 25; Q12: 17; Q13: 20;
%Q14: 25; reading and checking: 18 minutes. Total: 150 minutes.
%Existing questions: 64 points; new Q12--14: 12 each. Total: 100; grade=raw/10.
%These are planning estimates, not a timed student pilot.
%
%Question 9 follows the joint-distribution, moments and iterated-expectations
%material in Lecture 1 and Seminar 1. Question 10 follows standard errors,
%nonzero-null mean comparisons, critical-value tests, confidence intervals
%and sample-size planning in Lecture 2 and Seminar 2. Question 11 follows
%scalar OLS estimation, fitted values, residuals, fit, correlation and
%regressor-unit conversion in Lectures 1--3 and Seminars 1--3. The six
%written problems require numerical calculations rather than derivations.
%All numerical data are newly constructed for this exam; no multiple
%regression or later-course techniques are required.
%
%Production: use double-sided printing with long-edge binding. Copy quantity
%should equal enrolled students excluding cancelled exam sittings; do not
%request extra copies. Optical-reading instructions, if separately needed,
%belong in Moodle or on the classroom board, not in the student paper.
\end{document}
